\section{Analysis and limitations}
\label{sec:analysis}

\paragraph{Where \method fails.}
\label{sec:failures}
\method breaks \rOursBroken{} of the \rOrigCorrect{} originally correct clips (\steer: \rSteerBroken). In all of them the evaluation reader transcribes the original correctly while the pipeline's OCR reads it in at most 2\% of the frames (neon italics, extruded or condensed letters), so the preservation gate cannot fire; the outputs then show new text larger than its sign (4 clips), old text left visible next to the new text (5), wrong colour or contrast on lightboxes (2) or a misread stacked text (1). The more common failure does not affect readability: erasure leaves a visible patch, or the new glyphs are flat where the old ones were outlined, extruded or glowing. These cases explain the appearance distance in \cref{tab:main}.

\paragraph{The appearance gap to generative editing.}
\label{sec:limits}
Among the \rBothRepairN{} garbled clips that both methods repair, \steer's text is closer to the original in \rSteerCloserPct\%, so the gap is not an artefact of \steer leaving garbled text untouched. Four causes, in decreasing order of visibility: \method renders instead of generating, so it cannot reproduce extrusion, bevels, tube geometry or texture inside the strokes; its font library is small (\rZhFam{} Chinese families are ever selected, and Noto Sans SC alone renders \rNotoPct\% of the edited Chinese clips); its layout inherits the axis-aligned orientation of the OCR quad, so on tilted carriers the text can look skewed; and LaMa's fill is visible as a patch on textured carriers. Two observations suggest that the gap is closable without giving up correctness: the garbled source already contains a median of \rCharsEn\% of the target's letters (Chinese: \rCharsZh\%), so reusing them would extend ``change only what is wrong'' from words to characters once letter segmentation is reliable; and the two methods fail on different clips (at least one repairs \rAtLeastOne\%), which the verified selection of \cref{sec:exp_hybrid} exploits.

\paragraph{Other limitations.} The benchmark comes from one family of generators and asks for short sign texts; long texts are out of scope. Readability is measured by VLM readers, which favour large, high-contrast text and can transcribe a garbled original as the intended word; the second reader confirms the advantage over \steer with a smaller margin, and the third finds the methods level on English once its invented additions are discounted, so we report all of them. The appearance distance is glyph-agnostic and noisy at the clip level; it supports large differences such as the gap to \steer and the effect of SAT, not small ones, and it rewards a constant colour, which is why a flat paste scores well on it. The SAT study used four raters; the blind comparison of Banksy-V with \steer is in progress. The frozen front end is bit-exact within one software environment but not across OpenCV and CUDA versions, where its font matcher and a few OCR decisions diverge (about 3 readability points and 1.4\,dB of background PSNR on 71 clips); every comparison in this paper is therefore made within one environment and one evaluation run (supplement, Sec.~G). Outside repair, on editing correct text, the frozen system is clearly behind \steer (\cref{sec:exp_external}).
